\section{Kokkuvõte} \label{kokkuvote}

Käesolev bakalaureusetöö uuris tehisintellekti rakendamise teostatavust eestikeelse tehnilise informaatika lõputöö genereerimisel. Uurimuse käigus võrreldi nelja juhtivat suurt keelemudelit -- Google Gemini 3.1 Pro, OpenAI GPT-5.4, Anthropic Claude Sonnet 4.6 ja Anthropic Claude Opus 4.6 -- nende võimekuse osas luua struktureeritud, akadeemiliselt korrektset ja sisuliselt väärtuslikku eestikeelset teksti. Opus 4.6 puhul kasutati lisaks Claude Code'i agentset tööriista.

Uurimuse peamised tulemused on järgmised.

Esiteks, AI mudelid suudavad genereerida rahuldava kuni hea kvaliteediga eestikeelset akadeemilist teksti, eriti teksti struktureerimise ja kavandamise osas. Kontekstipõhine genereerimine, kus mudelitele anti ette Tartu Ülikooli juhendid ja näidistööd, parandas tulemusi keskmiselt 15--20\% võrreldes nullkonteksti genereerimisega.

Teiseks, nelja mudeli võrdlusest selgus, et Claude Opus 4.6 saavutas kõigis eksperimentides kõrgeima keskmise hinde (4,05--4,47 skaalal 5-st), ületades Sonnet 4.6 tulemust keskmiselt 0,35 palli võrra. Claude Sonnet 4.6 saavutas teise koha (3,86--4,17), olles eriti tugev akadeemilise stiili ja juhendite järgimise poolest ning pakkudes parimat hinna-kvaliteedi suhet. GPT-5.4 näitas kõige ühtlasemat eesti keele grammatilist kvaliteeti lühikeste tekstide puhul, samal ajal kui Gemini 3.1 Pro eeliseks oli märgatavalt suurem kontekstiaken. Mõtlemismudelite kasutamine parandas tulemusi 15--20\%, kuid suurendas kulusid 1,5--3 korda.

Kolmandaks, tervikliku lõputöö genereerimise eksperiment (eksperiment 6) paljastas mudelite vahel põhimõttelisi erinevusi. Claude Opus 4.6 genereeris Claude Code'i kaudu kõige mahukama väljundi -- 61-leheküljelise eestikeelse lõputöö paroolihalduri teemal, mis oli tehniliselt sügav ja strukturaalselt terviklik. Claude Sonnet 4.6 genereeris 40-leheküljelise eestikeelse lõputöö korrektsel teemal. GPT-5.4 keeldus terviklikku tööd kirjutamast, viidates akadeemilise aususe piirangutele, ja pakkus selle asemel 8-leheküljelise ingliskeelse teostatavusjuhendi. Gemini 3.1 Pro genereeris 20-leheküljelise eestikeelse töö, kuid vale teema kohta. Need tulemused näitavad, et mudelite juhiste järgimise võimekus, keelevalik ja teemale vastavus erinevad drastiliselt.

Neljandaks, kõige tõsisem probleem oli viidete hallutsineerimine -- mudelid leiutasid olematud allikaid 15--35\% juhtudest (Opus 4.6 ca 15\%, Sonnet 4.6 ca 20\%, GPT-5.4 ca 30\%, Gemini 3.1 Pro ca 35\%). See tähendab, et iga AI-genereeritud viide vajab inimese poolt kontrollimist. Samuti jäi AI-genereeritud teksti sisuline sügavus ja originaalsus oodatust madalamaks.

Viiendaks, AI kasutamise majanduslik kulu on Sonnet-klassi mudelitega mõistlik: \$23--42 API kaudu või \$20--40 kuutellimuste kaudu. Opus 4.6 on oluliselt kallim (\$158--200 API kaudu), kuid pakub kõrgeimat kvaliteeti. Mõlemad variandid on soodsamad kui professionaalse toimetamise teenus.

Uurimuse põhjal saab järeldada, et AI mudeleid saab efektiivselt kasutada lõputöö kirjutamise abivahendina, kuid mitte asendajana. AI tugevused -- struktureerimine, mustandi loomine, keeleline toimetamine ja LaTeX-vormindamine -- säästab üliõpilasele märkimisväärset aega. Samas on inimese panus endiselt asendamatu sisulise analüüsi, originaalsete järelduste, viidete kontrollimise ja valdkonnaspetsiifilise terminoloogia osas.

Soovitame kasutada Claude Opus 4.6 mudelit koos Claude Code'i agentse tööriistaga kõrgeima kvaliteedi saavutamiseks, eriti terviklike peatükkide ja keeruliste akadeemiliste tekstide genereerimisel. Piiratud eelarve korral on Claude Sonnet 4.6 parim hinna-kvaliteedi suhtega valik. GPT-5.4 sobib hästi keelelise korrektuuri jaoks ja lühikeste tekstilõikude genereerimiseks, kuid tervikliku töö genereerimisel tuleb arvestada mudeli ohutuspiirangutega. Gemini 3.1 Pro eeliseks on suur kontekstiaken mahuka konteksti töötlemiseks, kuid juhiste järgimist tuleb hoolikalt kontrollida. Ideaalne on Claude Code'i kasutamine koos Opus või Sonnet mudeliga (vt peatükk~\ref{claudecode}).

Käesolev töö panustab eestikeelsesse AI-uurimusse, pakkudes esimese süstemaatilise võrdluse kaasaegsete keelemudelite suutlikkusest eestikeelse akadeemilise teksti genereerimisel. Tervikliku lõputöö genereerimise eksperiment on eriti informatiivne, kuna paljastab mudelite käitumise erinevusi, mida lühikeste tekstilõikude testimine ei avalda. Lisaks annab töö praktilised soovitused Claude Code'i agentse tööriista kasutamiseks lõputöö genereerimisel (vt peatükk~\ref{claudecode}). Töö tulemused on kasulikud nii üliõpilastele, kes soovivad AI-d vastutustundlikult kasutada, kui ka õppejõududele ja ülikoolidele, kes kujundavad AI kasutamise poliitikat akadeemilises kontekstis.

Edaspidine uurimistöö võiks keskenduda spetsialiseeritud akadeemiliste AI tööriistade hindamisele, eesti keelele kohandatud mudelite (nt EstGPT) võrdlemisele suurte universaalsete mudelitega, RAG-süsteemide mõju hindamisele viidete kvaliteedile ning AI kasutamise mõju hindamisele üliõpilaste akadeemilisele arengule pikas perspektiivis.
